\newpage
\section*{\deva{परिशिष्ट} (Appendix): Glossary of Divine Terms}
\addcontentsline{toc}{section}{\deva{परिशिष्ट} (Appendix) — Glossary of Divine Terms}

Throughout the Hanuman Chalisa, Tulsidas references several deep cosmological and esoteric lists from Hindu theology. To avoid interrupting the flow of recitation, the detailed definitions of these profound terms are provided below.

\vspace{1em}
\subsection*{\textbf{1. The Eight Siddhis (Ashta-Siddhi)}}
In the 31st Chaupai, Mother Sita blesses Hanuman with the ability to grant the \textit{Ashta-Siddhi} (the eight supernatural Yogic powers). A master of these Siddhis can manipulate the fundamental building blocks of matter and space.

\begin{table}[h!]
\centering
\renewcommand{\arraystretch}{1.3}
\begin{tabularx}{\textwidth}{@{} l l X @{}}
\toprule
\textbf{Siddhi} & \textbf{Sanskrit} & \textbf{Supernatural Ability} \\
\midrule
\textbf{Anima} & \deva{अणिमा} & The power to shrink one's body to the size of an atom or sub-atomic particle (used to secretly enter Lanka). \\
\textbf{Mahima} & \deva{महिमा} & The power to expand one's body to an infinitely large, gargantuan size (used to terrify the demoness Surasa). \\
\textbf{Garima} & \deva{गरिमा} & The power to become infinitely heavy (used when Bhima could not lift even Hanuman's tail). \\
\textbf{Laghima} & \deva{लघिमा} & The power to become almost weightless or lighter than air (used for levitation and swift flight across the ocean). \\
\textbf{Prapti} & \deva{प्राप्ति} & The power to instantly acquire or reach anything anywhere, regardless of physical barriers. \\
\textbf{Prakamya} & \deva{प्राकाम्य} & The power to realize any desire, achieve irresistible willpower, and instantly materialize objects. \\
\textbf{Ishitva} & \deva{ईशित्व} & The power of absolute lordship and supreme dominance over nature and living beings. \\
\textbf{Vashitva} & \deva{वशित्व} & The ultimate power to conquer, subjugate, or mesmerize any entity to follow one's will. \\
\bottomrule
\end{tabularx}
\end{table}

\vspace{1em}
\subsection*{\textbf{2. The Nine Nidhis (Nava-Nidhi)}}
Alongside the 8 Siddhis, Sita also granted Hanuman dominion over the \textit{Nava-Nidhi} (the nine divine, inexhaustible treasures of Kubera, the god of wealth). These are not merely financial riches, but profound spiritual and cosmic abundances.

\begin{table}[h!]
\centering
\renewcommand{\arraystretch}{1.3}
\begin{tabularx}{\textwidth}{@{} l l X @{}}
\toprule
\textbf{Nidhi} & \textbf{Sanskrit} & \textbf{Divine Treasure} \\
\midrule
\textbf{Mahapadma} & \deva{महापद्म} & The ''Great Lotus Casket'' containing immense wealth, generating generations of prosperity and philanthropy. \\
\textbf{Padma} & \deva{पद्म} & The ''Lotus Casket'' offering abundant material riches, sattvic tendencies, and untainted gifts. \\
\textbf{Shankha} & \deva{शङ्ख} & The ''Conch Shell'' bringing self-sustaining wealth, typically enjoyed and isolated by the owner. \\
\textbf{Makara} & \deva{मकर} & The ''Crocodile'' granting martial proficiency, strategic intellect, and dominance in warfare/arms. \\
\textbf{Kacchapa} & \deva{कच्छप} & The ''Tortoise'' establishing wealth accumulated slowly, hidden secretly, and passed down carefully. \\
\textbf{Mukunda} & \deva{मुकुन्द} & The ''Quicksilver'' or ''Kettle'' bringing an aesthetic inclination towards fine arts, music, and royal patronage. \\
\textbf{Kunda} & \deva{कुन्द} & The ''Jasmine'' representing wealth related to morality, spiritual upliftment, and charitable endowments. \\
\textbf{Nila} & \deva{नील} & The ''Sapphire'' giving massive resources related to trade, commerce, and material exchange. \\
\textbf{Kharva} & \deva{खर्व} & The ''Dwarf'' representing wealth intertwined with mixed (both auspicious and complex) desires. \\
\bottomrule
\end{tabularx}
\end{table}

\vspace{1em}
\subsection*{\textbf{3. The Four Yugas (Chatur-Yuga)}}
In the 29th Chaupai, Tulsidas proclaims that Hanuman's glory is famous across all four ages (\textit{Yugas}). According to Vedic cosmology, time is cyclical, divided into four grand planetary epochs.

\begin{table}[h!]
\centering
\renewcommand{\arraystretch}{1.3}
\begin{tabularx}{\textwidth}{@{} l l X @{}}
\toprule
\textbf{Yuga} & \textbf{Sanskrit} & \textbf{Era Characteristics} \\
\midrule
\textbf{Satya Yuga} & \deva{सत्ययुग} & The Age of Truth and absolute Purity. Humanity is deeply spiritual, peaceful, and perfectly completely aligned with Dharma. \\
\textbf{Treta Yuga} & \deva{त्रेतायुग} & The Silver Age. Dharma slightly declines. This is the era in which Lord Rama and Hanuman incarnated to establish righteousness. \\
\textbf{Dvapara Yuga} & \deva{द्वापरयुग} & The Bronze Age. Morality drops to 50\%. This era featured the Mahabharata war and the incarnation of Lord Krishna. \\
\textbf{Kali Yuga} & \deva{कलियुग} & The Iron Age (our current era). An age of deep ignorance, conflict, and spiritual amnesia. Hanuman remains immortal (\textit{Chiranjeevi}) physically present on Earth specifically to guide seekers through this dark age. \\
\bottomrule
\end{tabularx}
\end{table}

\vspace{1em}
\subsection*{\textbf{4. The Four Fruits of Life (Phala Chaari)}}
In the very first opening Doha, Tulsidas prays for the ''four fruits'' (\textit{Phala Chaari}). In Vedic philosophy, these are the \textit{Purusharthas}—the four fundamental grand goals or aims of human life that lead to complete fulfillment.

\begin{table}[h!]
\centering
\renewcommand{\arraystretch}{1.3}
\begin{tabularx}{\textwidth}{@{} l l X @{}}
\toprule
\textbf{Fruit} & \textbf{Sanskrit} & \textbf{Life Goal} \\
\midrule
\textbf{Dharma} & \deva{धर्म} & Righteousness, moral duty, and living in cosmic law and harmony. The foundational aim of a noble life. \\
\textbf{Artha} & \deva{अर्थ} & Material prosperity, wealth, economic security, and physical well-being required to sustain life and family. \\
\textbf{Kama} & \deva{काम} & Desire, pleasure, emotional fulfillment, aesthetics, and the enjoyment of the world in a balanced way. \\
\textbf{Moksha} & \deva{मोक्ष} & The ultimate liberation. Freedom from the cycle of birth and death (\textit{Samsara}) and union with the divine. \\
\bottomrule
\end{tabularx}
\end{table}
